\section{Findings that came later}
\label{sec:afterwards}

The sections above were written before the revision was carried out. Four
things were found while carrying it out, each of which changed what was done.
They are recorded here rather than folded into the earlier sections, so that
the order in which the evidence arrived stays legible.

\subsection{The FK Monte-Carlo does not measure the FK diagram}

Section~\ref{sec:sampling} treats the appendix Monte-Carlo as a check that
could not have caught the interpolation error, because it reads the same
tabulated cumulant. That is true, but it understates the situation. With the
corrected vertex the Monte-Carlo points depart from the analytic channel by a
factor two at $2'$ and carry the wrong sign beyond $6'$, which sent us to the
estimator itself.

Its expectation can be computed exactly at the diagram's order, by pure linear
algebra with no realisations, and doing so reproduces all $21$ probe points to
about $1.5$ seed standard errors. It decomposes into two terms, neither of
which is the FK diagram. First, the variance-reduced estimator injects the
third cumulant only on the observable-side leg of the vertex, so of the three
placements that make up $\zeta = Q\!:\!VV + \text{perms}$ it captures one; in
the eigenbasis of the node covariance the captured fraction is
$d_q d_r/(d_p d_q + d_q d_r + d_r d_p)$, which kernel-weighted comes to
$0.245$ at $0.5'$, passes through zero near $10'$ and is negative beyond.
Second, a finite-$\sigma_\lambda$ term enters because the AR(1) smearing
decoheres the node-local vertex from the smeared covariance, releasing the
cancellation of the tabulated tensors' $\lambda$-grid graininess; smoothing
both tensors makes the model $\sigma_\lambda$-flat at exactly the share
prediction, which is the attribution.

The historical agreement at $\sigma_\lambda=8$ was these two errors summing
through unity at a value that had itself been chosen so that the ratio came
out one. The FK amplitude also scales close to linearly with
$\sigma_\lambda$ ($8.70$, $19.8$ and $37.9\times10^{-6}$ at $0.5'$ for
$\sigma_\lambda=4,8,16$), with an intercept consistent with zero, so no range
of that estimator is a converged measurement of FK.

The FK points were therefore withdrawn from the figure. The channel is now
checked by a direct quadrature contraction of the tabulated cumulant with the
FK kernel,
\[
  \mathrm{FK}(\gamma) = 2\!\int\!\mathrm{d}v\, W(v)\,H(v)\,
     \Bigl[-\!\sum_c \zeta^{(6)}_{cc B}(\gamma; v)\Bigr],
  \qquad
  H(v) = D(v)^4\!\!\int_v^{\lambda_s}\!\!\frac{W(u)}{D(u)^4}\,\mathrm{d}u ,
\]
which shares only the vertex table with the workflow and reproduces the folded
channel to $0.99$--$1.02$ over $0.5'$--$5'$ and $1.07$ at $17.3'$. The FF
channel keeps its Monte-Carlo points: FF is a term in the equation of motion
and can be differenced at a fixed random draw, so its noise scales with the
effect itself, whereas a third moment measured from realisations scatters
according to the Gaussian fluctuations of the field however small the cumulant
is.

\subsection{The multi-redshift FF slices had never been computed}

The apparent low-$\ell$ excess at $z_s=1.7$ in the multi-redshift figure was
not physics and not the interpolation error either. The cache filled the FF
slices at $z_s<5$ with the exact $z_s=5$ result scaled by the \emph{per-%
$\gamma$} ratio $f_K(z)/f_K(5)$. At large separation $f_K$ sits on its
numerical floor, so that ratio is noise over noise: the tails carried
$10^{-7}$-level dips and went negative at $z_s=1$, which the FF moment cannot
do, being a connected part plus $\langle\kappa\rangle^2>0$. Through the
truncated-range transform this became an excess of order the Order-0 signal.

Five genuine per-redshift sweeps were run in its place, one single-threaded
process per source distance, about $3.2$ hours each. Two checks: each slice's
large-separation plateau matches the deterministic $\langle\kappa\rangle^2(z_s)$,
obtained by propagating the expectation of the same recursions with no
Monte-Carlo, to $1.022$, $1.005$, $1.005$, $0.999$ and $1.008$ at
$z_s=1,\,1.7,\,2.5,\,3.2,\,4$; and a $z_s=5$ control, run the same way with
fresh caches, reproduces the production FF \emph{bit-identically} at all forty
separations. The production path is a real computation, deterministically
reproducible, and this is the same path.

\subsection{The multipole sum stops converging beyond about a degree}

Separation enters the cumulant only through $P_{\ell_3}(\cos\gamma)$. That
factor is unity for every term at coincidence, so the sum is coherent there
and each doubling of the cutoff adds to it. Once $\ell_3\gamma\gtrsim1$ it
instead oscillates with period $\Delta\ell_3=2\pi/\gamma$ and is damped as
$(\ell_3\gamma)^{-1/2}$, so the high-multipole terms alternate in sign and
largely cancel. The low-multipole branch, where $\ell_3\gamma\ll1$, stays
almost constant: measured on the $z=5.7$ shell, $\zeta_{TTT}$ for $\ell\le60$
is $+1.33\times10^{-16}$ at every separation, while the $\ell>60$ part falls
from $-3.7\times10^{-13}$ at $0.5'$ to $-1.5\times10^{-16}$ at $232'$. The two
necessarily cross, and near the crossing $\zeta$ is a small difference of two
comparable numbers computed by different methods.

Three independent symptoms follow, and all three place the boundary near a
degree. Two evaluations of the same multipole sum, one wide window against
eight disjoint ones, agree to $3$--$6\%$ below $85'$ and differ by factors of
several to several tens beyond $117'$. Two independent evaluations of
$\xi_{\rm FK}$, the quadrature above against the workflow fold, agree to
$0.4$--$6\%$ below $20'$ and differ by a factor $5.2$ at $114'$. And raising
the cutoff from $7680$ to $15360$ changes $\xi_{\rm FK}$ by a uniform
$1.08$--$1.14$ for $\gamma\le56'$ but by $0.69$, $0.28$, $2.59$ and $0.73$ at
$114'$, $294'$, $597'$ and $1535'$.

The consequence for the harmonic-space numbers is direct, because low
multipoles are dominated by the large-separation part of the two-point
function. Zeroing $\xi_{\rm FK}$ beyond a cut and re-transforming moves the
$\ell=2$ correction between $-10.9\%$ and $+2.7\%$ as the cut is varied, while
$\ell=100$ stays at $1.14$--$1.24\%$ throughout. FK is therefore quoted only
for $\gamma\lesssim1^\circ$ and $\ell\gtrsim50$, and drawn faint outside that
range.

\subsection{The $\zeta$-slices figure was still at the old cutoff}

It was the last product drawn at $\ell_{\max}=1000$ while everything else had
moved to $15360$, and the cutoff changes the panels rather than only their
normalisation: at the $z=3.1$ shell the peak $|\zeta_{TTT}|$ grows by $13.7$
and the normalised value at $5'$ moves from $-0.87$ to $-0.27$, so the printed
panels showed a cumulant falling about three times more slowly in angle than
the one the numbers come from. Rebuilt at the converged cutoff by appending
four flat-sky Born-Limber windows to $(8000,15360]$, the exact Wigner-3j ladder
staying at $\ell\le60$, and clipped to $\gamma\le85'$ for the reason above.
Its caption also carried a blanket ``change sign with angular scale'' which is
not true at the converged cutoff: within the plotted range $\zeta_{TTP}$ never
changes sign, $\zeta_{TTT}$ and $\zeta_{B}$ change sign at $59'$ only for
$z\gtrsim4$, and $\zeta_{D}$ from $z\simeq1$.
